\section{环境}
\subsection{用户}
id命令查看用户信息，用户定义在\texttt{/etc/passwd}文件中，属组定义在\texttt{/etc/group}文件中，
一般来说，会为每个用户创建一个与用户同名的单成员属组。如果想切换用户可以使用su命令，可以使用sudo来以超级用户权限执行命令。
编辑用户可以使用useradd、userdel和usermod命令。

\subsection{软件包安装}
在从源码安装软件时，通常需要先运行 ./configure。
configure 是源码中提供的一个 Shell 脚本，它主要用于：分析和检测当前系统的构建环境、
根据环境情况对源代码和编译选项进行适当调整、检查是否已安装编译所需的外部依赖和工具。
在完成这些工作后，configure 会生成一个适合本机环境的 Makefile 文件。接着，就可以通过 make 命令来编译源码。
编译后通常可以通过make install来安装软件，通常安装在/usr/local/bin下。

\subsection{环境变量}
Shell保存了两种数据类型：Shell变量和环境变量。
可以使用printenv命令查看环境变量、set查看所有变量、alias查看别名。
\begin{table}[H]
	\centering
	\begin{tabular}{|l|l|}
		\hline
		\textbf{变量} & \textbf{说明}                \\
		\hline
		DISPLAY     & 屏幕名称，通常为:0，表示X服务器生成的第一个屏幕。 \\
		\hline
		HOME        & 当前用户的主目录。                  \\
		\hline
		LANG        & 系统语言和区域设置。                 \\
		\hline
		OLDPWD      & 上一个工作目录。                   \\
		\hline
		PATH        & 由冒号分隔的目录列表，用于可执行文件的搜索。     \\
		\hline
		PWD         & 当前工作目录。                    \\
		\hline
		SHELL       & 当前用户的默认Shell。              \\
		\hline
		TERM        & 终端类型。                      \\
		\hline
		USER        & 当前用户名。                     \\
		\hline
	\end{tabular}
	\caption{常见环境变量}
\end{table}

\subsection{配置文件}
Shell 会话通常分为两类：登录会话和非登录会话，二者在启动时所读取的配置文件有所不同。
当修改了配置文件后，可以通过 source 或\verb|.|命令使改动立即生效。
登录会话在启动时需要用户输入用户名和密码进行身份验证，而非登录会话则无需凭证即可进入。
通常，通过图形界面（GUI）启动的终端属于非登录会话。
\begin{table}[H]
	\centering
	\begin{tabular}{|l|l|}
		\hline
		\textbf{文件}            & \textbf{说明}                           \\
		\hline
		\verb|/etc/profile|    & 应用于所有用户的全局配置。                         \\
		\hline
		\verb|~/.bash_profile| & 用户个人的登录 Shell 配置文件。                   \\
		\hline
		\verb|~/.bash_login|   & 如果 \verb|~/.bash_profile| 不存在，则读取该文件。 \\
		\hline
		\verb|~/.profile|      & 如果前两者都不存在，则读取该文件。                     \\
		\hline
	\end{tabular}
	\caption{登录 Shell 读取的配置文件}
\end{table}

\begin{table}[H]
	\centering
	\begin{tabular}{|l|p{6cm}|}
		\hline
		\textbf{文件}             & \textbf{说明}                                           \\
		\hline
		\verb|/etc/bash.bashrc| & 应用于所有用户的全局配置。                                         \\
		\hline
		\verb|~/.bashrc|        & 用户个人的非登录 Shell 配置文件，通常会在 \verb|~/.bash_profile| 中被调用。 \\
		\hline
	\end{tabular}
	\caption{非登录 Shell 读取的配置文件}
\end{table}
